\documentclass[11pt, a4paper]{article}

% ── Paquetes ──────────────────────────────────────────────────────────────────
\usepackage[utf8]{inputenc}
\usepackage[T1]{fontenc}
\usepackage[spanish]{babel}
\usepackage{geometry}
\geometry{margin=2.5cm}
\usepackage{graphicx}
\usepackage{xcolor}
\usepackage{enumitem}
\usepackage{booktabs}
\usepackage{tabularx}
\usepackage{listings}
\usepackage{amsmath}
\usepackage{hyperref}
\usepackage{fancyhdr}
\usepackage{titlesec}
\usepackage{tcolorbox}
\usepackage{fontawesome5}
\usepackage{tikz}
\usepackage{array}
\usetikzlibrary{arrows.meta, positioning, shapes.geometric}

% ── Colores ───────────────────────────────────────────────────────────────────
\definecolor{primary}{HTML}{1A237E}
\definecolor{accent}{HTML}{E65100}
\definecolor{success}{HTML}{2E7D32}
\definecolor{danger}{HTML}{C62828}
\definecolor{warning}{HTML}{EF6C00}
\definecolor{info}{HTML}{1565C0}
\definecolor{codebg}{HTML}{F5F5F5}
\definecolor{codeframe}{HTML}{BDBDBD}
\definecolor{get}{HTML}{2E7D32}
\definecolor{post}{HTML}{1565C0}
\definecolor{patch}{HTML}{EF6C00}

% ── Configuración de listings ─────────────────────────────────────────────────
\lstset{
  basicstyle=\ttfamily\small,
  backgroundcolor=\color{codebg},
  frame=single,
  rulecolor=\color{codeframe},
  breaklines=true,
  showstringspaces=false,
  extendedchars=true,
  inputencoding=utf8,
  literate=
    {á}{{\'a}}{1} {é}{{\'e}}{1} {í}{{\'i}}{1} {ó}{{\'o}}{1} {ú}{{\'u}}{1}
    {ñ}{{\~n}}{1} {Á}{{\'A}}{1} {É}{{\'E}}{1} {Í}{{\'I}}{1} {Ó}{{\'O}}{1}
    {Ú}{{\'U}}{1} {Ñ}{{\~N}}{1}
}

% ── Cajas (solo para texto, NO para código) ───────────────────────────────────
\newtcolorbox{infobox}[1][]{
  colback=info!5, colframe=info!80!black,
  fonttitle=\bfseries, title={\faInfoCircle\quad Nota}, #1
}

\newtcolorbox{warningbox}[1][]{
  colback=warning!5, colframe=warning!80!black,
  fonttitle=\bfseries, title={\faExclamationCircle\quad Importante}, #1
}

% ── Comando para método HTTP ──────────────────────────────────────────────────
\newcommand{\httpmethod}[2]{%
  \colorbox{#1}{\textcolor{white}{\textbf{\small\textsf{#2}}}}%
}

% ── Comando para encabezado de endpoint ───────────────────────────────────────
\newcommand{\routeheader}[3]{%
  \noindent\rule{\textwidth}{1.5pt}\\[0.3em]
  \httpmethod{#1}{#2}\;\texttt{\bfseries #3}\\[0.3em]
  \rule{\textwidth}{0.4pt}
  \vspace{0.3cm}
}

% ── Comando para encabezado de contrato ───────────────────────────────────────
\newcommand{\contractheader}[1]{%
  \noindent\colorbox{primary!10}{%
    \parbox{\dimexpr\textwidth-2\fboxsep}{%
      \textbf{\faFileCode\quad #1}%
    }%
  }
  \vspace{0.3cm}
}

% ── Encabezado y pie de página ────────────────────────────────────────────────
\pagestyle{fancy}
\fancyhf{}
\fancyhead[L]{\small\textcolor{primary}{\textbf{Grupo 6} --- Contrato API Despacho y Logística}}
\fancyhead[R]{\small\textcolor{primary}{v1.1}}
\fancyfoot[C]{\small\textcolor{gray}{\thepage}}
\renewcommand{\headrulewidth}{0.4pt}

% ── Formato de títulos ────────────────────────────────────────────────────────
\titleformat{\section}
  {\Large\bfseries\color{primary}}{\thesection.}{0.5em}{}[\titlerule]
\titleformat{\subsection}
  {\large\bfseries\color{primary!80}}{\thesubsection}{0.5em}{}
\titleformat{\subsubsection}
  {\normalsize\bfseries\color{primary!60}}{\thesubsubsection}{0.5em}{}

\hypersetup{
  colorlinks=true, linkcolor=primary, urlcolor=info, citecolor=accent,
}

% ══════════════════════════════════════════════════════════════════════════════
\begin{document}

% ── Portada ───────────────────────────────────────────────────────────────────
\begin{titlepage}
\centering
\vspace*{2.5cm}

{\Huge\bfseries\textcolor{primary}{Contrato API}}\\[0.3cm]
{\LARGE\textcolor{primary!70}{Servicio de Despacho y Logística}}\\[0.5cm]
{\Large\textcolor{accent}{\textbf{Grupo 6}}}\\[2cm]

{\large Especificación de Endpoints REST}\\[0.2cm]
{\large Contratos JSON de Eventos}\\[0.2cm]
{\large Matriz de Dependencias}\\[3cm]

\begin{tabular}{rl}
  \textbf{Proyecto:}     & Marketplace en la Nube \\
  \textbf{Versión API:}  & \texttt{v1.1} \\
  \textbf{Base URL:}     & \texttt{https://g6-despacho.onrender.com/api/v1} \\
  \textbf{Fecha:}        & 17 de Junio de 2026 \\
\end{tabular}

\vfill



\end{titlepage}

\tableofcontents
\newpage


% ══════════════════════════════════════════════════════════════════════════════
\part{Contrato REST API}
% ══════════════════════════════════════════════════════════════════════════════

\section{Responsabilidad del Servicio y Dominio}

\subsection{Objetivo}

El servicio de \textbf{Despacho y Logística (G6)} es responsable de gestionar el ciclo de vida completo de los envíos dentro del ecosistema del Marketplace en la Nube.

\subsection{Qué resuelve este servicio}

\begin{itemize}[noitemsep]
  \item Crear despachos a partir de pedidos confirmados por G5.
  \item Asignar un identificador único (\texttt{shipment\_id}) a cada envío.
  \item Gestionar la máquina de estados del envío (\texttt{PENDING} $\rightarrow$ \texttt{IN\_TRANSIT} $\rightarrow$ \texttt{DELIVERED}, etc.).
  \item Publicar eventos de cambio de estado para que G8 notifique al usuario y G7 genere reportes.
  \item Exponer endpoints REST para consulta de tracking (G1) y listados (G7).
  \item Almacenar y administrar la tabla \texttt{shipments} como fuente de verdad.
\end{itemize}

\subsection{Qué queda fuera de alcance}

\begin{itemize}[noitemsep]
  \item \textbf{Gestión de pedidos:} responsabilidad de G5 (Order Management).
  \item \textbf{Procesamiento de pagos:} responsabilidad de G8 (Pago Simulado).
  \item \textbf{Generación de reportes:} responsabilidad de G7 (Reportería y Batch).
  \item \textbf{Interfaz de usuario:} responsabilidad de G1 (Frontend).
  \item \textbf{Autenticación y autorización:} delegada a G1/G8.
\end{itemize}

\subsection{Dueño del Dato (Ownership)}

\begin{table}[h!]
\centering
\renewcommand{\arraystretch}{1.3}
\begin{tabularx}{\textwidth}{l l l X}
  \toprule
  \textbf{Dato} & \textbf{Dueño} & \textbf{Acceso G6} & \textbf{Observación} \\
  \midrule
  \texttt{shipments}        & G6 (propio)  & Lectura/Escritura & Fuente de verdad del envío. \\
  \texttt{order\_id}        & G5 (ajeno)   & Solo lectura      & Recibido en el POST, nunca modificado. \\
  \texttt{customer\_name}   & G5 (ajeno)   & Solo lectura      & Copiado del pedido para el envío. \\
  \texttt{payment\_status}  & G8 (ajeno)   & No accede          & G6 no consulta ni modifica pagos. \\
  \texttt{report\_data}     & G7 (ajeno)   & No accede          & G7 consume datos vía REST/eventos. \\
  \bottomrule
\end{tabularx}
\caption{Ownership de datos: qué datos son propios y cuáles son consultados.}
\end{table}


\section{Información General}
\label{sec:general-info}

\begin{table}[h!]
\centering
\renewcommand{\arraystretch}{1.4}
\begin{tabularx}{\textwidth}{l X}
  \toprule
  \textbf{Campo} & \textbf{Valor} \\
  \midrule
  Título          & API de Despacho y Logística --- Grupo 6 \\
  Versión         & 1.0.0 \\
  Base URL        & \texttt{https://g6-despacho.onrender.com/api/v1} \\
  Formato         & JSON (\texttt{application/json}) \\
  Autenticación   & Ninguna (proyecto académico) \\
  Encoding        & UTF-8 \\
  \bottomrule
\end{tabularx}
\end{table}

\subsection{Convenciones}

\begin{itemize}[noitemsep]
  \item Todos los endpoints requieren los siguientes headers obligatorios:
  \begin{itemize}
    \item \texttt{X-Request-Id}: uuid
    \item \texttt{X-Correlation-Id}: uuid
    \item \texttt{X-Consumer}: group-name
    \item \texttt{Idempotency-Key}: uuid
  \end{itemize}
  \item Todos los endpoints retornan \texttt{application/json}.
  \item Las fechas usan formato ISO 8601: \texttt{YYYY-MM-DDTHH:mm:ssZ}.
  \item Los IDs de shipment usan formato \texttt{SHP-<uuid>} (generados por nuestro servicio).
  \item Los errores siguen estructura estándar (ver Sección~\ref{sec:errores}).
\end{itemize}

\subsection{Resumen de Endpoints}

\begin{table}[h!]
\centering
\renewcommand{\arraystretch}{1.5}
\begin{tabularx}{\textwidth}{l l X l}
  \toprule
  \textbf{Método} & \textbf{Ruta} & \textbf{Descripción} & \textbf{Consumidor} \\
  \midrule
  \httpmethod{post}{POST}   & \texttt{/shipments}                  & Crear un nuevo despacho          & G5 \\
  \httpmethod{get}{GET}     & \texttt{/shipments/\{id\}}           & Consultar estado de un despacho  & G1, G5 \\
  \httpmethod{get}{GET}     & \texttt{/shipments?order\_id=\{x\}}  & Buscar despacho por pedido       & G5 \\
  \httpmethod{patch}{PATCH} & \texttt{/shipments/\{id\}}           & Actualizar estado del despacho   & Interno \\
  \httpmethod{get}{GET}     & \texttt{/shipments}                  & Listar todos los despachos       & G7 \\
  \httpmethod{get}{GET}     & \texttt{/health}                     & Health check del servicio        & Todos \\
  \bottomrule
\end{tabularx}
\caption{Resumen de endpoints expuestos por el Grupo 6.}
\end{table}


\section{Modelo de Datos: Shipment}
\label{sec:modelo}

\begin{table}[h!]
\centering
\renewcommand{\arraystretch}{1.3}
\begin{tabularx}{\textwidth}{l l l c X}
  \toprule
  \textbf{Campo} & \textbf{Tipo} & \textbf{Ejemplo} & \textbf{Req.} & \textbf{Descripción} \\
  \midrule
  \texttt{shipment\_id}        & string   & \texttt{SHP-a1b2c3d4}        & ---  & ID único generado por G6. \\
  \texttt{order\_id}           & string   & \texttt{ORD-20260611-001}     & Sí   & ID del pedido (de G5). \\
  \texttt{customer\_name}      & string   & \texttt{María G.}             & Sí   & Nombre del destinatario. \\
  \texttt{address}             & string   & \texttt{Av. Prov. 1234}       & Sí   & Dirección de entrega. \\
  \texttt{city}                & string   & \texttt{Santiago}             & Sí   & Ciudad de destino. \\
  \texttt{weight\_kg}          & number   & \texttt{3.2}                  & Sí   & Peso del paquete en kg. \\
  \texttt{status}              & string   & \texttt{PENDING}              & ---  & Estado actual. \\
  \texttt{created\_at}         & datetime & \texttt{2026-06-11T14:00Z}    & ---  & Fecha de creación. \\
  \texttt{updated\_at}         & datetime & \texttt{2026-06-11T15:30Z}    & ---  & Última actualización. \\
  \texttt{estimated\_delivery} & datetime & \texttt{2026-06-14T18:00Z}    & ---  & Fecha estimada. \\
  \bottomrule
\end{tabularx}
\caption{Esquema del recurso \texttt{Shipment}. ``Req.'' indica si es requerido en el \texttt{POST}.}
\end{table}


\subsection{Estados válidos (Máquina de Estados de Doble Nivel)}
\label{sec:estados}

Para evitar incongruencias visuales en el checkout y tracking, el sistema divide el estado físico del comercial.

\subsubsection{5.1 Nivel Físico (Por Caja - G6)}
Estados almacenados en la tabla \texttt{shipments}.

\begin{center}
\begin{tikzpicture}[
  state/.style={rectangle, rounded corners=5pt, draw=primary, thick,
    fill=primary!8, minimum width=2.6cm, minimum height=0.9cm,
    font=\small\bfseries, text=primary},
  error/.style={rectangle, rounded corners=5pt, draw=danger, thick,
    fill=danger!8, minimum width=2.6cm, minimum height=0.9cm,
    font=\small\bfseries, text=danger},
  arrow/.style={-{Stealth[length=3mm]}, thick, primary!70},
  errarrow/.style={-{Stealth[length=3mm]}, thick, danger!70, dashed},
]
  \node[state]                     (pending)   {PENDING};
  \node[state, right=1.2cm of pending]   (transit)   {IN\_TRANSIT};
  \node[state, right=1.2cm of transit]   (delivered) {DELIVERED};
  \node[error, below=1.2cm of pending]   (cancelled) {CANCELLED};
  \node[error, below=1.2cm of transit]   (failed)    {FAILED};
  \node[error, below=1.2cm of delivered] (returned)  {RETURNED};
  \draw[arrow]    (pending)   -- (transit);
  \draw[arrow]    (transit)   -- (delivered);
  \draw[errarrow] (pending)   -- (cancelled);
  \draw[errarrow] (transit)   -- (failed);
  \draw[errarrow] (delivered) -- (returned);
\end{tikzpicture}
\end{center}

\begin{table}[h!]
\centering
\renewcommand{\arraystretch}{1.3}
\begin{tabularx}{\textwidth}{l l l X}
  \toprule
  \textbf{Estado} & \textbf{Valor} & \textbf{Tipo} & \textbf{Transiciones válidas} \\
  \midrule
  Pendiente   & \texttt{PENDING}     & Inicial & $\rightarrow$ \texttt{IN\_TRANSIT}, \texttt{CANCELLED} \\
  En tránsito & \texttt{IN\_TRANSIT} & Normal  & $\rightarrow$ \texttt{DELIVERED}, \texttt{FAILED} \\
  Entregado   & \texttt{DELIVERED}   & Final   & $\rightarrow$ \texttt{RETURNED} \\
  Cancelado   & \texttt{CANCELLED}   & Final   & Ninguna (terminal). \\
  Fallido     & \texttt{FAILED}      & Final   & Ninguna (terminal). \\
  Devuelto    & \texttt{RETURNED}    & Final   & Ninguna (terminal). \\
  \bottomrule
\end{tabularx}
\caption{Tabla de transiciones de estado a nivel físico (por caja).}
\end{table}

\subsubsection{5.2 Nivel Comercial / Global (Por Pedido - G1/G5)}
Estado visualizado por el cliente, calculado al vuelo leyendo todas las cajas de un \texttt{order\_id}:
\begin{itemize}[noitemsep]
  \item \textbf{PENDING:} Todas las cajas están pendientes.
  \item \textbf{PARTIALLY\_DELIVERED:} Al menos una caja fue entregada, el resto sigue en tránsito.
  \item \textbf{DELIVERED:} 100\% de las cajas entregadas.
  \item \textbf{HAS\_ISSUES:} Al menos un paquete registró fallo o devolución.
\end{itemize}




\section{Lógica de Negocio: Tarifa Única Consolidada}

El endpoint \texttt{/quotes} debe calcular la suma total de las cajas para retornar un único cobro.

\textbf{Paso 1: Peso Facturable (Factor de Conversión: 4000)}
Para cada paquete, se cobra por el mayor valor entre el peso físico y el volumétrico:
$$ \text{Peso Volumétrico} = \frac{\text{Largo (cm)} \times \text{Ancho (cm)} \times \text{Alto (cm)}}{4000} $$
$$ \text{Peso Facturable} = \max(\text{Peso Físico}, \text{Peso Volumétrico}) $$

\textbf{Paso 2: Tarifario Simulado por Macrozona}
\begin{itemize}[noitemsep]
  \item \textbf{Misma Zona} (Ej. CD Centro a Santiago): Base \$3000 + (\$500 x Kg Facturable).
  \item \textbf{Zona Adyacente} (Ej. CD Centro a Norte/Sur): Base \$5000 + (\$800 x Kg Facturable).
  \item \textbf{Zona Extrema} (Ej. CD Norte a Sur): Base \$8000 + (\$1200 x Kg Facturable).
\end{itemize}

\textbf{Paso 3: Ecuación Total Consolidada}
$$ \text{Costo Total} = \sum_{i=1}^{n} \left( \text{Costo Base}_i + (\text{Peso Facturable}_i \times \text{Valor Kg}_i) \right) $$

% ══════════════════════════════════════════════════════════════════════════════
\section{Endpoints Detallados}
% ══════════════════════════════════════════════════════════════════════════════

% ── POST /shipments ───────────────────────────────────────────────────────────
\subsection{Crear Despacho}

\routeheader{post}{POST}{/api/v1/shipments}

\textbf{Descripción:} Crea una nueva orden de despacho. Llamado por
\textbf{G5 (Pedidos o Order Management)} cuando un pedido ha sido confirmado y pagado.

\vspace{0.3cm}
\textbf{Request Body} (\texttt{application/json}):

\begin{lstlisting}
{
  "order_id": "ORD-20260611-001",
  "customer_name": "María González",
  "address": "Av. Providencia 1234, Depto 56",
  "city": "Santiago",
  \"packages\": [
    { \"origin_cd\": \"NORTE\", \"weight_kg\": 2.5, \"dimensions_cm\": { \"length\": 40, \"width\": 30, \"height\": 20 } },
    { \"origin_cd\": \"CENTRO\", \"weight_kg\": 1.0, \"dimensions_cm\": { \"length\": 15, \"width\": 10, \"height\": 5 } }
  ]
}
\end{lstlisting}

\textbf{Campos requeridos:}
\begin{table}[h!]
\small
\begin{tabularx}{\textwidth}{l l l X}
  \toprule
  \textbf{Campo} & \textbf{Tipo} & \textbf{Req.} & \textbf{Validaciones} \\
  \midrule
  \texttt{order\_id}      & string & Sí & No vacío. Único. \\
  \texttt{customer\_name} & string & Sí & No vacío. Máx. 200 chars. \\
  \texttt{address}        & string & Sí & No vacío. Máx. 500 chars. \\
  \texttt{city}           & string & Sí & No vacío. Máx. 100 chars. \\
  \texttt{weight\_kg}     & number & Sí & Mayor a 0. Máx. 500 kg. \\
  \bottomrule
\end{tabularx}
\end{table}

\textbf{Response --- 201 Created:}

\begin{lstlisting}
[
  \"SHP-1\",
  \"SHP-2\"
]
\end{lstlisting}

\textbf{Errores:}
\begin{table}[h!]
\small
\begin{tabularx}{\textwidth}{l l X}
  \toprule
  \textbf{Código} & \textbf{Nombre} & \textbf{Causa} \\
  \midrule
  400 & Bad Request          & JSON mal formado o campos faltantes. \\
  409 & Conflict             & Ya existe despacho para ese \texttt{order\_id}. \\
  422 & Unprocessable Entity & Datos inválidos (ej: peso negativo). \\
  \bottomrule
\end{tabularx}
\end{table}

\noindent\rule{\textwidth}{0.4pt}


\noindent\rule{\textwidth}{0.4pt}

% ── POST /shipments/quotes ───────────────────────────────────────────────────
\subsection{Cotizar Envíos}

\routeheader{post}{POST}{/api/v1/shipments/quotes}

\textbf{Descripción:} Cotiza envíos según peso volumétrico y orígenes. Llamado por \textbf{G4 (Inventario y Checkout)}.
Fórmula: \texttt{Peso Volumétrico (kg) = (Largo\_cm * Ancho\_cm * Alto\_cm) / 4000}.

\vspace{0.3cm}
\textbf{Request Body} (\texttt{application/json}):

\begin{lstlisting}
{
  "city": "Santiago",
  "address": "Av. Providencia 1234",
  "packages": [
    {
      "origin_cd": "NORTE",
      "weight_kg": 2.5,
      "dimensions_cm": { "length": 40, "width": 30, "height": 20 }
    }
  ]
}
\end{lstlisting}

\textbf{Response --- 200 OK:}

\begin{lstlisting}
{
  "total_shipping_cost": 6600,
  "currency": "CLP"
}
\end{lstlisting}

% ── GET /shipments/{id} ──────────────────────────────────────────────────────
\subsection{Consultar Despacho por ID}

\routeheader{get}{GET}{/api/v1/shipments/\{shipment\_id\}}

\textbf{Descripción:} Retorna estado actual y datos completos de un despacho.
Llamado por \textbf{G1 (Frontend)} y \textbf{G5 (Pedidos o Order Management)}.

\vspace{0.3cm}
\textbf{Response --- 200 OK:}

\begin{lstlisting}
[
  {
  "shipment_id": "SHP-a1b2c3d4",
  "order_id": "ORD-20260611-001",
  "customer_name": "María González",
  "address": "Av. Providencia 1234, Depto 56",
  "city": "Santiago",
  "weight_kg": 3.2,
  "status": "IN_TRANSIT",
  "created_at": "2026-06-11T14:00:00Z",
  "updated_at": "2026-06-11T14:30:00Z",
  "estimated_delivery": "2026-06-14T18:00:00Z"
}
]
\end{lstlisting}

\textbf{Errores:}
\begin{table}[h!]
\small
\begin{tabularx}{\textwidth}{l l X}
  \toprule
  \textbf{Código} & \textbf{Nombre} & \textbf{Causa} \\
  \midrule
  404 & Not Found & No existe despacho con ese ID. \\
  \bottomrule
\end{tabularx}
\end{table}

\noindent\rule{\textwidth}{0.4pt}

% ── GET /shipments?order_id= ─────────────────────────────────────────────────
\subsection{Buscar Despacho por Order ID}

\routeheader{get}{GET}{/api/v1/shipments?order\_id=\{order\_id\}}

\textbf{Descripción:} Busca un despacho por el ID de pedido de G5.

\vspace{0.3cm}
\textbf{Response --- 200 OK:}

\begin{lstlisting}
[
  {
  "shipment_id": "SHP-a1b2c3d4",
  "order_id": "ORD-20260611-001",
  "status": "PENDING",
  "created_at": "2026-06-11T14:00:00Z",
  "estimated_delivery": "2026-06-14T18:00:00Z"
}
]
\end{lstlisting}

\textbf{Errores:}
\begin{table}[h!]
\small
\begin{tabularx}{\textwidth}{l l X}
  \toprule
  \textbf{Código} & \textbf{Nombre} & \textbf{Causa} \\
  \midrule
  400 & Bad Request & Falta parámetro \texttt{order\_id}. \\
  404 & Not Found   & No existe despacho para ese pedido. \\
  \bottomrule
\end{tabularx}
\end{table}

\noindent\rule{\textwidth}{0.4pt}

% ── PATCH /shipments/{id} ────────────────────────────────────────────────────
\subsection{Actualizar Estado del Despacho}

\routeheader{patch}{PATCH}{/api/v1/shipments/\{shipment\_id\}}

\textbf{Descripción:} Actualiza el estado de un despacho. Uso \textbf{interno}
(simulación automática) o por \textbf{G5} para cancelar.

\vspace{0.3cm}
\textbf{Request Body:}

\begin{lstlisting}
{
  "status": "IN_TRANSIT"
}
\end{lstlisting}

\textbf{Response --- 200 OK:}

\begin{lstlisting}
{
  "shipment_id": "SHP-a1b2c3d4",
  "order_id": "ORD-20260611-001",
  "status": "IN_TRANSIT",
  "previous_status": "PENDING",
  "updated_at": "2026-06-11T14:30:00Z"
}
\end{lstlisting}

\begin{infobox}
  Al cambiar el estado exitosamente, el servicio emite automáticamente
  el evento correspondiente hacia G8 (ver Parte~II).
\end{infobox}

\textbf{Errores:}
\begin{table}[h!]
\small
\begin{tabularx}{\textwidth}{l l X}
  \toprule
  \textbf{Código} & \textbf{Nombre} & \textbf{Causa} \\
  \midrule
  400 & Bad Request & Estado no es un valor válido. \\
  404 & Not Found   & No existe despacho con ese ID. \\
  409 & Conflict    & Transición no permitida (ej: \texttt{DELIVERED} $\rightarrow$ \texttt{PENDING}). \\
  \bottomrule
\end{tabularx}
\end{table}

\noindent\rule{\textwidth}{0.4pt}

% ── GET /shipments (lista) ───────────────────────────────────────────────────
\subsection{Listar Despachos}

\routeheader{get}{GET}{/api/v1/shipments}

\textbf{Descripción:} Lista de despachos con filtros. Para \textbf{G7 (Reportería, Batch y Streaming)}.

\vspace{0.3cm}
\textbf{Query Parameters (opcionales):}
\begin{table}[h!]
\small
\begin{tabularx}{\textwidth}{l l l X}
  \toprule
  \textbf{Param} & \textbf{Tipo} & \textbf{Default} & \textbf{Descripción} \\
  \midrule
  \texttt{status} & string & (todos) & Filtrar por estado. \\
  \texttt{limit}  & int    & 50      & Máx. resultados. \\
  \texttt{offset} & int    & 0       & Paginación. \\
  \bottomrule
\end{tabularx}
\end{table}

\textbf{Response --- 200 OK:}

\begin{lstlisting}
{
  "total": 142,
  "limit": 50,
  "offset": 0,
  "shipments": [
    {
      "shipment_id": "SHP-a1b2c3d4",
      "order_id": "ORD-20260611-001",
      "status": "DELIVERED",
      "city": "Santiago",
      "created_at": "2026-06-11T14:00:00Z",
      "updated_at": "2026-06-11T16:00:00Z"
    }
  ]
}
\end{lstlisting}

\noindent\rule{\textwidth}{0.4pt}

% ── GET /health ──────────────────────────────────────────────────────────────
\subsection{Health Check}

\routeheader{get}{GET}{/api/v1/health}

\begin{lstlisting}
{
  "status": "ok",
  "service": "g6-despacho",
  "version": "1.0.0",
  "timestamp": "2026-06-11T14:00:00Z"
}
\end{lstlisting}

\noindent\rule{\textwidth}{0.4pt}


% ── Estructura Estándar de Errores ───────────────────────────────────────────
\section{Estructura Estándar de Errores}
\label{sec:errores}

Todos los errores siguen el mismo formato base. Opcionalmente, se puede incluir un campo \texttt{"details"} para proveer más contexto de validación:

\begin{lstlisting}
{
  "timestamp": "2026-06-11T14:05:00Z",
  "status": 422,
  "code": "VALIDATION_ERROR",
  "message": "El campo weight_kg debe ser mayor a 0.",
  "correlationId": "req-123e4567"
}
\end{lstlisting}

\begin{table}[h!]
\centering
\renewcommand{\arraystretch}{1.3}
\begin{tabularx}{\textwidth}{l l X}
  \toprule
  \textbf{status} & \textbf{code} & \textbf{Uso} \\
  \midrule
  400 & \texttt{BAD\_REQUEST}      & JSON malformado, campos faltantes. \\
  404 & \texttt{NOT\_FOUND}        & Recurso no encontrado. \\
  409 & \texttt{CONFLICT}          & Duplicado o transición inválida. \\
  422 & \texttt{VALIDATION\_ERROR} & Datos que no pasan validación. \\
  500 & \texttt{INTERNAL\_ERROR}   & Error inesperado del servidor. \\
  \bottomrule
\end{tabularx}
\caption{Códigos de error estandarizados.}
\end{table}


% ══════════════════════════════════════════════════════════════════════════════
\newpage
\part{Contrato de Eventos JSON}
% ══════════════════════════════════════════════════════════════════════════════

\section{Mecanismo de Comunicación (Pub/Sub)}

La arquitectura de eventos se basa en un modelo \textbf{Pub/Sub} (Publish/Subscribe) utilizando un \textbf{Event Broker} (ej. Upstash Kafka o Supabase Realtime). 
Esto permite desacoplar los servicios y procesar los eventos de forma 100\% asíncrona.

\subsection{Flujo de Eventos}

\begin{enumerate}[noitemsep]
  \item Un despacho cambia de estado en nuestro servicio (G6).
  \item G6 publica un evento JSON (Publish) en el tópico global \texttt{shipment-events} del Broker.
  \item Los servicios consumidores (G8, G7) que están suscritos (Subscribe) a este tópico reciben el mensaje asíncronamente.
  \item El Broker se encarga de asegurar la entrega, eliminando el acoplamiento directo entre G6 y los consumidores.
\end{enumerate}


\section{Formato Estándar del Evento}
\label{sec:formato-evento}

Todos los eventos emitidos por G6 siguen esta estructura. Cada vez que un despacho cambie de estado, el servicio del Grupo 6 generará y publicará el siguiente payload JSON en el tópico global \texttt{shipment-events}. Los servicios consumidores (Grupo 8 y Grupo 7) deberán estar suscritos a dicho tópico para procesar la información de forma asíncrona:

\contractheader{Envelope del Evento (estructura común)}

\begin{lstlisting}
{
  "eventId": "evt-550e8400-e29b-41d4-a716-446655440000",
  "eventType": "ShipmentDelivered",
  "version": "1.0",
  "occurredAt": "2026-06-11T15:00:00Z",
  "producer": "g6-despacho",
  "correlationId": "req-123e4567",
  "payload": {
    \"package_index\": 1,
    \"total_packages\": 2, ... }
}
\end{lstlisting}

\begin{table}[h!]
\centering
\renewcommand{\arraystretch}{1.3}
\begin{tabularx}{\textwidth}{l l X}
  \toprule
  \textbf{Campo} & \textbf{Tipo} & \textbf{Descripción} \\
  \midrule
  \texttt{eventId}       & string (UUID) & Identificador único del evento. \\
  \texttt{eventType}     & string        & Tipo del evento (ver catálogo). \\
  \texttt{version}       & string        & Versión del esquema del evento. \\
  \texttt{occurredAt}    & datetime      & Momento del evento (ISO 8601). \\
  \texttt{producer}      & string        & Siempre \texttt{"g6-despacho"}. \\
  \texttt{correlationId} & string        & ID de correlación para trazar flujo. \\
  \texttt{payload}       & object        & Datos específicos del evento. \\
  \bottomrule
\end{tabularx}
\caption{Campos del envelope de evento.}
\end{table}


\section{Catálogo de Eventos}

% ── ShipmentCreated ──────────────────────────────────────────────────────────
\subsection{ShipmentCreated}

Se emite cuando G5 crea exitosamente un nuevo despacho.

\contractheader{ShipmentCreated}

\begin{lstlisting}
{
  "eventId": "evt-550e8400-e29b-41d4-a716-446655440001",
  "eventType": "ShipmentCreated",
  "version": "1.0",
  "occurredAt": "2026-06-11T14:00:00Z",
  "producer": "g6-despacho",
  "correlationId": "req-123e4567",
  "payload": {
    \"package_index\": 1,
    \"total_packages\": 2,
    "shipment_id": "SHP-a1b2c3d4",
    "order_id": "ORD-20260611-001",
    "customer_name": "María González",
    "address": "Av. Providencia 1234, Depto 56",
    "city": "Santiago",
    "weight_kg": 3.2,
    "status": "PENDING",
    "estimated_delivery": "2026-06-14T18:00:00Z"
  }
}
\end{lstlisting}

\textbf{Consumidores:} G8 (notifica al cliente), G7 (registro).

\vspace{0.4cm}

% ── ShipmentInTransit ────────────────────────────────────────────────────────
\subsection{ShipmentInTransit}

Se emite cuando el paquete sale a reparto.

\contractheader{ShipmentInTransit}

\begin{lstlisting}
{
  "eventId": "evt-550e8400-e29b-41d4-a716-446655440002",
  "eventType": "ShipmentInTransit",
  "version": "1.0",
  "occurredAt": "2026-06-11T14:30:00Z",
  "producer": "g6-despacho",
  "correlationId": "req-123e4567",
  "payload": {
    \"package_index\": 1,
    \"total_packages\": 2,
    "shipment_id": "SHP-a1b2c3d4",
    "order_id": "ORD-20260611-001",
    "previous_status": "PENDING",
    "new_status": "IN_TRANSIT",
    "customer_name": "María González",
    "city": "Santiago"
  }
}
\end{lstlisting}

\textbf{Consumidores:} G8 (``Tu paquete va en camino'').

\vspace{0.4cm}

% ── ShipmentDelivered ────────────────────────────────────────────────────────
\subsection{ShipmentDelivered}

Se emite cuando el paquete es entregado exitosamente.

\contractheader{ShipmentDelivered}

\begin{lstlisting}
{
  "eventId": "evt-550e8400-e29b-41d4-a716-446655440003",
  "eventType": "ShipmentDelivered",
  "version": "1.0",
  "occurredAt": "2026-06-11T15:00:00Z",
  "producer": "g6-despacho",
  "correlationId": "req-123e4567",
  "payload": {
    \"package_index\": 1,
    \"total_packages\": 2,
    "shipment_id": "SHP-a1b2c3d4",
    "order_id": "ORD-20260611-001",
    "previous_status": "IN_TRANSIT",
    "new_status": "DELIVERED",
    "customer_name": "María González",
    "city": "Santiago",
    "delivered_at": "2026-06-11T15:00:00Z"
  }
}
\end{lstlisting}

\textbf{Consumidores:} G8 (``Tu paquete ha llegado''), G7 (métricas).

\vspace{0.4cm}

% ── ShipmentCancelled ────────────────────────────────────────────────────────
\subsection{ShipmentCancelled}

Se emite cuando un despacho es cancelado antes de salir a reparto.

\contractheader{ShipmentCancelled}

\begin{lstlisting}
{
  "eventId": "evt-550e8400-e29b-41d4-a716-446655440004",
  "eventType": "ShipmentCancelled",
  "version": "1.0",
  "occurredAt": "2026-06-11T14:10:00Z",
  "producer": "g6-despacho",
  "correlationId": "req-123e4567",
  "payload": {
    \"package_index\": 1,
    \"total_packages\": 2,
    "shipment_id": "SHP-a1b2c3d4",
    "order_id": "ORD-20260611-001",
    "previous_status": "PENDING",
    "new_status": "CANCELLED",
    "customer_name": "María González",
    "reason": "Pedido anulado por el cliente"
  }
}
\end{lstlisting}

\textbf{Consumidores:} G8 (notifica cancelación), G5 (confirma).

\vspace{0.4cm}

% ── ShipmentFailed ───────────────────────────────────────────────────────────
\subsection{ShipmentFailed}

Se emite cuando un intento de entrega falla.

\contractheader{ShipmentFailed}

\begin{lstlisting}
{
  "eventId": "evt-550e8400-e29b-41d4-a716-446655440005",
  "eventType": "ShipmentFailed",
  "version": "1.0",
  "occurredAt": "2026-06-11T16:00:00Z",
  "producer": "g6-despacho",
  "correlationId": "req-123e4567",
  "payload": {
    \"package_index\": 1,
    \"total_packages\": 2,
    "shipment_id": "SHP-a1b2c3d4",
    "order_id": "ORD-20260611-001",
    "previous_status": "IN_TRANSIT",
    "new_status": "FAILED",
    "customer_name": "María González",
    "reason": "Destinatario ausente"
  }
}
\end{lstlisting}

\textbf{Consumidores:} G8 (``Hubo un problema con tu entrega''), G5.

\vspace{0.4cm}

% ── ShipmentReturned ─────────────────────────────────────────────────────────
\subsection{ShipmentReturned}

Se emite cuando un paquete entregado es devuelto.

\contractheader{ShipmentReturned}

\begin{lstlisting}
{
  "eventId": "evt-550e8400-e29b-41d4-a716-446655440006",
  "eventType": "ShipmentReturned",
  "version": "1.0",
  "occurredAt": "2026-06-12T10:00:00Z",
  "producer": "g6-despacho",
  "correlationId": "req-123e4567",
  "payload": {
    \"package_index\": 1,
    \"total_packages\": 2,
    "shipment_id": "SHP-a1b2c3d4",
    "order_id": "ORD-20260611-001",
    "previous_status": "DELIVERED",
    "new_status": "RETURNED",
    "customer_name": "María González",
    "reason": "Producto incorrecto"
  }
}
\end{lstlisting}

\textbf{Consumidores:} G8, G5 (gestionar reembolso), G7 (métricas).


% ══════════════════════════════════════════════════════════════════════════════
\newpage
\part{Contratos Inter-Grupo}
% ══════════════════════════════════════════════════════════════════════════════

\section{Contrato con G5 (Pedidos o Order Management) --- Entrada}
\label{sec:contrato-g5}

Este contrato define \textbf{qué datos nos envía G5} cuando un pedido es
confirmado y necesita ser despachado.

\subsection{Flujo de Integración}

\begin{enumerate}[noitemsep]
  \item El cliente completa el checkout y el pago es procesado por G8.
  \item G5 (Pedidos o Order Management) consolida el pedido y nos notifica para despacho.
  \item G5 hace una solicitud HTTP POST a nuestro \texttt{/api/v1/shipments}.
  \item Respondemos con el \texttt{shipment\_id} y estado \texttt{PENDING}.
  \item G5 almacena el \texttt{shipment\_id} para tracking y consultas futuras.
\end{enumerate}

\subsection{JSON que G5 nos envía}

\contractheader{Request de G5 hacia G6}

\begin{lstlisting}
{
  "order_id": "ORD-20260611-001",
  "customer_name": "María González",
  "address": "Av. Providencia 1234, Depto 56",
  "city": "Santiago",
  "weight_kg": 3.2
}
\end{lstlisting}

\subsection{JSON que G6 le responde a G5}

\contractheader{Response de G6 hacia G5}

\begin{lstlisting}
[
  \"SHP-1\",
  \"SHP-2\"
]
\end{lstlisting}

\subsection{Preguntas pendientes para G5}

\begin{warningbox}[title={\faExclamationCircle\quad Confirmar con G5}]
\begin{enumerate}[noitemsep]
  \item Formato de \texttt{order\_id}: propuesto como string libre.
  \item ¿Necesitan campos adicionales en la respuesta?
  \item ¿Nos envían el peso o lo calculamos nosotros?
  \item ¿Cancelan vía \texttt{PATCH} o necesitan endpoint dedicado?
\end{enumerate}
\end{warningbox}


\section{Contrato con G8 (Pago Simulado y Notificaciones) --- Salida}
\label{sec:contrato-g8}

Define \textbf{qué datos le enviamos a G8} en cada cambio de estado.

\subsection{Flujo de Integración}

\begin{enumerate}[noitemsep]
  \item Un shipment cambia de estado (vía PATCH o simulación).
  \item Construimos el evento JSON correspondiente.
  \item Publicamos el evento en el tópico \texttt{shipment-events} del Event Broker.
  \item G8 consume el evento desde el tópico y notifica al usuario final.
\end{enumerate}

\subsection{JSON que G6 publica en el Event Broker (para G8)}

\contractheader{Evento publicado por G6 (ejemplo: entrega)}

\begin{lstlisting}
{
  "eventId": "evt-550e8400-e29b-41d4-a716-446655440003",
  "eventType": "ShipmentDelivered",
  "version": "1.0",
  "occurredAt": "2026-06-11T15:00:00Z",
  "producer": "g6-despacho",
  "correlationId": "req-123e4567",
  "payload": {
    \"package_index\": 1,
    \"total_packages\": 2,
    "shipment_id": "SHP-a1b2c3d4",
    "order_id": "ORD-20260611-001",
    "previous_status": "IN_TRANSIT",
    "new_status": "DELIVERED",
    "customer_name": "María González",
    "city": "Santiago",
    "delivered_at": "2026-06-11T15:00:00Z"
  }
}
\end{lstlisting}

\subsection{Comportamiento esperado de G8 (Pago Simulado y Notificaciones)}

\begin{itemize}[noitemsep]
  \item \textbf{Procesamiento:} El servicio del Grupo 8 consumirá el evento
    desde el tópico \texttt{shipment-events} de forma asíncrona. G6 no recibe
    ninguna respuesta directa ni código HTTP de parte de G8.
  \item \textbf{Acuse de recibo (ACK):} Una vez procesada la información
    localmente por su lógica interna, el componente consumidor de G8 emitirá un
    \textbf{ACK} (\textit{Acknowledgment}) automático hacia el Event Broker para
    confirmar la recepción exitosa y permitir la liberación del mensaje en la
    cola. Este ACK es una señal de bajo nivel gestionada por la librería del
    Broker, no un payload JSON personalizado.
  \item \textbf{Acción final:} El Grupo 8 será responsable de gatillar la
    alerta o correo correspondiente al usuario final utilizando los datos del
    payload transmitido (\texttt{order\_id}, \texttt{status},
    \texttt{customer\_name}).
\end{itemize}

\begin{infobox}
  En una arquitectura Pub/Sub real, el productor (G6) nunca ve el ACK del
  consumidor (G8). El ACK viaja exclusivamente entre el consumidor y el
  Broker, y su único efecto es marcar el mensaje como procesado en la cola.
  Esto garantiza el desacoplamiento total entre ambos servicios.
\end{infobox}

\subsection{Preguntas pendientes para G8}

\begin{warningbox}[title={\faExclamationCircle\quad Confirmar con G8}]
\begin{enumerate}[noitemsep]
  \item ¿Campos adicionales necesarios? (email, teléfono)
  \item ¿Qué eventos les interesan? ¿Los 6 o solo algunos?
  \item Coordinar acceso y credenciales al Event Broker.
\end{enumerate}
\end{warningbox}


\section{Contrato con G7 (Reportería, Batch y Streaming) --- Salida}

G7 consume datos para dashboards de rendimiento.

\subsection{Opciones de Integración}

\begin{enumerate}[noitemsep]
  \item \textbf{Vía API REST:} \texttt{GET /shipments} con filtros.
  \item \textbf{Vía Eventos:} Suscripción al tópico \texttt{shipment-events} en el Broker.
\end{enumerate}

\subsection{Datos relevantes para G7}

\begin{table}[h!]
\centering
\renewcommand{\arraystretch}{1.3}
\begin{tabularx}{\textwidth}{l X}
  \toprule
  \textbf{Métrica} & \textbf{Cómo la obtienen} \\
  \midrule
  Total despachos       & \texttt{GET /shipments} $\rightarrow$ \texttt{total}. \\
  Por estado            & Filtrar con \texttt{?status=DELIVERED}. \\
  Tasa de éxito         & \texttt{DELIVERED} / total. \\
  Tasa de fallos        & (\texttt{FAILED} + \texttt{RETURNED}) / total. \\
  Tiempo de entrega     & \texttt{updated\_at} - \texttt{created\_at} en \texttt{DELIVERED}. \\
  \bottomrule
\end{tabularx}
\end{table}


\section{Contrato con G1 (Frontend) --- Consulta}

G1 consume nuestro endpoint de tracking.

\textbf{Endpoint:} \texttt{GET /api/v1/shipments/\{shipment\_id\}}

\begin{table}[h!]
\centering
\renewcommand{\arraystretch}{1.3}
\begin{tabularx}{\textwidth}{l X}
  \toprule
  \textbf{Campo} & \textbf{Uso en la UI} \\
  \midrule
  \texttt{status}              & Paso actual en la barra de progreso. \\
  \texttt{estimated\_delivery} & ``Llega el 14 de Junio''. \\
  \texttt{city}                & ``Tu paquete va camino a Santiago''. \\
  \texttt{updated\_at}         & ``Última actualización: hace 30 min''. \\
  \bottomrule
\end{tabularx}
\end{table}


% ══════════════════════════════════════════════════════════════════════════════
\newpage
\part{Matriz de Dependencias y Arquitectura}
% ══════════════════════════════════════════════════════════════════════════════

\section{Matriz de Dependencias}

\begin{table}[h!]
\centering
\renewcommand{\arraystretch}{1.5}
\begin{tabularx}{\textwidth}{c c l l X}
  \toprule
  \textbf{Origen} & \textbf{Destino} & \textbf{Tipo} & \textbf{Protocolo} & \textbf{Datos} \\
  \midrule
  G5  & G6  & Síncrono  & REST POST    & Datos del pedido. \\
  G5  & G6  & Síncrono  & REST GET     & Consulta por \texttt{order\_id}. \\
  G5  & G6  & Síncrono  & REST PATCH   & Cancelación. \\
  G1  & G6  & Síncrono  & REST GET     & Tracking por \texttt{shipment\_id}. \\
  G6     & Broker (Kafka/Supabase) & Asíncrono & Publish (Pub)   & Mensaje en tópico \texttt{shipment-events} \\
  Broker & G8 (Pago Simulado y Notificaciones)     & Asíncrono & Subscribe (Sub) & Consume desde tópico \texttt{shipment-events} \\
  Broker & G7 (Reportería, Batch y Streaming)        & Asíncrono & Subscribe (Sub) & Consume desde tópico \texttt{shipment-events} \\
  G7 & G6  & Síncrono  & REST GET     & Lista para reportes. \\
  \bottomrule
\end{tabularx}
\caption{Matriz completa de dependencias del Grupo 6.}
\end{table}


\section{Diagrama de Arquitectura (C4 --- Nivel 2: Contenedores)}

\begin{center}
{\small\textbf{System Container Diagram --- Ecosistema Logístico (perspectiva G6 Despachos)}}
\vspace{0.5cm}

\begin{tikzpicture}[
  box/.style={rectangle, rounded corners=8pt, draw=primary, thick, fill=primary!6,
    minimum width=3.4cm, minimum height=1.8cm, align=center,
    font=\small\bfseries, text=primary},
  extbox/.style={rectangle, rounded corners=8pt, draw=info, thick, fill=info!6,
    minimum width=3.4cm, minimum height=1.8cm, align=center,
    font=\small\bfseries, text=info},
  db/.style={cylinder, draw=accent, thick, fill=accent!8, shape border rotate=90,
    minimum width=2.4cm, minimum height=1.4cm, aspect=0.3,
    font=\small\bfseries, text=accent},
  brokerbox/.style={rectangle, rounded corners=8pt, draw=warning, thick, fill=warning!5,
    minimum width=3.4cm, minimum height=1.8cm, align=center,
    font=\small\bfseries, text=warning!80!black},
  arr/.style={-{Stealth[length=3mm]}, thick},
  tech/.style={font=\tiny\itshape, text=gray},
]
  % ── Sistemas externos (izquierda) ──
  \node[extbox]                    (g5)  {G5 --- Pedidos\\\relax{\tiny[Python / FastAPI]}};
  \node[extbox, above=0.8cm of g5] (g4)  {G4 --- Checkout\\\relax{\tiny[Node.js]}};
  \node[extbox, above=0.8cm of g4] (g1)  {G1 --- Frontend\\\relax{\tiny[React / Next.js]}};

  % ── Contenedor principal (centro) ──
  \node[box, right=2.8cm of g5]    (g6)  {G6 --- Despachos\\\relax{\tiny[Python / FastAPI]}};
  \node[db, below=1.8cm of g6]     (db)  {Supabase};
  \node[tech, below=0.1cm of db]   {[PostgreSQL]};

  % ── Event Broker + consumidores (derecha) ──
  \node[brokerbox, right=2.8cm of g6] (broker) {Event Broker\\\relax{\tiny[Kafka / Supabase Realtime]}};
  \node[extbox, above=1.8cm of broker] (g8)  {G8 --- Pago Simulado\\y Notificaciones\\\relax{\tiny[Node.js]}};
  \node[extbox, below=1.8cm of broker] (g7)  {G7 --- Reportería,\\Batch y Streaming\\\relax{\tiny[Python / Spark]}};

  % ── Flechas REST (síncronas) ──
  \draw[arr, info] (g5)  -- node[above, font=\scriptsize] {\texttt{POST/GET/PATCH}} (g6);
  \draw[arr, info] (g4)  -- node[above, font=\scriptsize] {\texttt{POST /quotes}} (g6);
  \draw[arr, info] (g1)  -| node[above left, font=\scriptsize] {\texttt{GET /shipments}} (g6);
  \draw[arr, primary] (g6) -- node[left, font=\scriptsize] {SQL} (db);

  % ── Flechas Eventos (asíncronas) ──
  \draw[arr, success] (g6) -- node[above, font=\scriptsize] {Publish Event} (broker);
  \draw[arr, success] (broker) -- node[right, font=\scriptsize] {Subscribe} (g8);
  \draw[arr, success] (broker) -- node[right, font=\scriptsize] {Subscribe} (g7);
  \draw[arr, info, dashed] (g7.south) -- ++(0,-1.0) -| node[below, near start, font=\scriptsize] {\texttt{GET}} (g6.south);

\end{tikzpicture}
\end{center}

\vspace{0.5cm}
\begin{center}
{\footnotesize
  \textcolor{info}{\rule{0.6cm}{2pt}} Comunicación síncrona (REST/HTTP) \qquad
  \textcolor{success}{\rule{0.6cm}{2pt}} Comunicación asíncrona (Pub/Sub) \qquad
  \textcolor{info}{- - -} Consulta periódica
}
\end{center}


\section{Esquema SQL de la Base de Datos (PostgreSQL/Supabase)}

Para soportar el modelo Multi-Origen y asegurar la trazabilidad y resiliencia de eventos, se implementan tres tablas:

\subsection{Tabla Principal: shipments}
Se elimina la restricción UNIQUE en \texttt{order\_id} para permitir 1:N.

\begin{lstlisting}
CREATE TABLE shipments (
    shipment_id        TEXT PRIMARY KEY,
    order_id           TEXT NOT NULL,
    customer_name      TEXT NOT NULL,
    address            TEXT NOT NULL,
    city               TEXT NOT NULL,
    origin_cd          TEXT NOT NULL, -- 'NORTE', 'CENTRO', 'SUR'
    volumetric_weight  REAL NOT NULL,
    shipping_cost      INTEGER NOT NULL,
    weight_kg          REAL NOT NULL CHECK (weight_kg > 0),
    status             TEXT NOT NULL DEFAULT 'PENDING'
                       CHECK (status IN (
                           'PENDING', 'IN_TRANSIT', 'DELIVERED',
                           'CANCELLED', 'FAILED', 'RETURNED'
                       )),
    created_at         TIMESTAMP NOT NULL DEFAULT NOW(),
    updated_at         TIMESTAMP NOT NULL DEFAULT NOW(),
    estimated_delivery TIMESTAMP
);
\end{lstlisting}

\subsection{Tabla de Historial Temporal: shipment\_status\_history}
Garantiza la trazabilidad (línea de tiempo) inmutable para el frontend.

\begin{lstlisting}
CREATE TABLE shipment_status_history (
    id SERIAL PRIMARY KEY,
    shipment_id TEXT REFERENCES shipments(shipment_id),
    status TEXT NOT NULL,
    created_at TIMESTAMP NOT NULL DEFAULT NOW()
);
\end{lstlisting}

\subsection{Tabla de Resiliencia: outbox\_events (Patrón Outbox)}
Mitiga caídas del Event Broker.

\begin{lstlisting}
CREATE TABLE outbox_events (
    id SERIAL PRIMARY KEY,
    event_type TEXT NOT NULL,
    payload JSONB NOT NULL,
    status TEXT NOT NULL DEFAULT 'PENDING', -- PENDING, PUBLISHED
    created_at TIMESTAMP NOT NULL DEFAULT NOW()
);
\end{lstlisting}

\begin{infobox}
  Esquema para Supabase (PostgreSQL). Si se usa Firebase,
  se adaptará a documentos JSON con los mismos campos.
\end{infobox}


\section{Supuestos, Riesgos y Mitigaciones}

\subsection{Supuestos}

\begin{enumerate}[noitemsep]
  \item \textbf{Formato de entrada estandarizado:} G5 (Pedidos) envía los datos del pedido con todos los campos obligatorios definidos en el contrato REST (\texttt{order\_id}, \texttt{customer\_name}, \texttt{address}, \texttt{city}, \texttt{weight\_kg}).
  \item \textbf{Idempotencia del \texttt{order\_id}:} Cada pedido genera exactamente un envío. Si G5 reintenta el \texttt{POST}, el sistema rechaza el duplicado con \texttt{409 Conflict}.
  \item \textbf{Disponibilidad del Event Broker:} El Event Broker (Kafka o Supabase Realtime) está operativo y accesible. La publicación de eventos es \emph{at-least-once}.
  \item \textbf{Base de datos centralizada:} G6 es el único grupo que escribe en la tabla \texttt{shipments}. Otros grupos solo consultan vía REST.
  \item \textbf{Conectividad cloud:} Todos los servicios se despliegan en la nube y se comunican a través de HTTPS o protocolo de mensajería seguro.
  \item \textbf{Autenticación delegada:} La autenticación del usuario final la gestiona G1/G8. G6 confía en los headers \texttt{X-Consumer} y \texttt{X-Correlation-Id} para trazabilidad.
\end{enumerate}

\subsection{Riesgos y Mitigaciones}

\begin{table}[!ht]
\centering
\renewcommand{\arraystretch}{1.4}
\begin{tabularx}{\textwidth}{l X X l}
  \toprule
  \textbf{ID} & \textbf{Riesgo} & \textbf{Mitigación} & \textbf{Impacto} \\
  \midrule
  R1 & Caída del Event Broker impide que G8 notifique al usuario final. & Cola de reintentos con backoff exponencial; tabla \texttt{outbox} local para eventos pendientes. & Alto \\
  R2 & G5 envía datos incompletos o con formato inválido. & Validación estricta en el endpoint \texttt{POST /shipments} con retorno \texttt{422} y mensaje descriptivo. & Medio \\
  R3 & Duplicación de envíos por reintentos de G5 sin \texttt{Idempotency-Key}. & Verificación estricta mediante \texttt{Idempotency-Key} a nivel de base de datos. & Medio \\
  R4 & Latencia alta en la base de datos Supabase por volumen de consultas de G7. & Índices en \texttt{status} y \texttt{created\_at}; endpoint paginado con \texttt{limit/offset}. & Medio \\
  R5 & Transición de estado inválida (e.g., de \texttt{DELIVERED} a \texttt{PENDING}). & Máquina de estados con validación explícita de transiciones permitidas antes de aplicar el PATCH. & Bajo \\
  R6 & Pérdida de trazabilidad entre servicios ante fallos parciales. & Propagación obligatoria de \texttt{X-Correlation-Id} en todos los requests y eventos publicados. & Medio \\
  \bottomrule
\end{tabularx}
\caption{Riesgos identificados para el servicio de Despachos (G6).}
\end{table}

\subsection{Manejo de Errores}

Todos los errores retornados por la API siguen la estructura estándar definida en la Sección~\ref{sec:general-info} (Información General). Además:

\begin{itemize}[noitemsep]
  \item Los errores de validación (\texttt{422}) incluyen el campo específico que falló en el \texttt{message}.
  \item Los conflictos de idempotencia (\texttt{409}) retornan el \texttt{shipment\_id} existente para que el consumidor pueda recuperarse.
  \item Los errores de transición de estado inválida (\texttt{422}) indican el estado actual y los estados de destino válidos.
  \item Todos los errores incluyen el \texttt{correlationId} para facilitar la depuración entre grupos.
\end{itemize}


\section{Historial de Versiones}

\begin{table}[h!]
\centering
\begin{tabularx}{\textwidth}{l l l X}
  \toprule
  \textbf{Ver.} & \textbf{Fecha} & \textbf{Autor} & \textbf{Cambios} \\
  \midrule
  1.0 & 16 Jun 2026 & Grupo 6 & Versión inicial. \\
  1.1 & 17 Jun 2026 & Grupo 6 & Soporte Multi-Origen y Tarifas por Peso Volumétrico (v1.1). \\
  \bottomrule
\end{tabularx}
\end{table}


\end{document}
